\documentclass[12pt,a4paper]{article}
\usepackage[utf8]{inputenc}
\usepackage[spanish]{babel}
\usepackage{amsmath,amssymb,amsfonts}
\usepackage{geometry}
\usepackage{hyperref}
\geometry{margin=2.5cm}

\title{\textbf{Constantes Físicas en Acústica}}
\author{Compendio de Constantes Fundamentales}
\date{\today}

\begin{document}

\maketitle

\section*{Introducción}
La acústica es la rama de la física que estudia las ondas mecánicas en medios gaseosos, líquidos y sólidos, abarcando vibraciones, sonido, ultrasonido e infrasonido. Sus constantes de referencia definen la propagación ondulatoria y los umbrales de audición humana.

\section{Velocidad del Sonido en Aire Estándar Seco ($c_s$ o $v_s$)}

\subsection*{1. Significado, Símbolo y Explicación}
Simbolizada por $c_s$ (o $v_s$), es la velocidad de propagación de una onda de presión longitudinal a través del aire seco a la temperatura de referencia de $0^\circ\text{C}$ ($273.15\text{ K}$) y $1\text{ atm}$ de presión. Se rige por la ecuación de Laplace:
\begin{equation*}
c_s = \sqrt{\frac{\gamma R T}{M}}
\end{equation*}
donde $\gamma \approx 1.40$ es el índice adiabático del aire, $R$ la constante de los gases y $M$ la masa molar del aire seco.

\subsection*{2. Valores Típicos en Ingeniería y Ciencia}
A $20^\circ\text{C}$ (temperatura ambiente habitual en ingeniería acústica):
\begin{equation*}
c_s \approx 343 \text{ m/s} \quad \text{o} \quad 343.2 \text{ m/s}
\end{equation*}
A $0^\circ\text{C}$ (estándar ISO 2533):
\begin{equation*}
c_s \approx 331.3 \text{ m/s}
\end{equation*}

\subsection*{3. Decimales Completos y Valor Estándar ISO}
Según la Atmósfera Estándar Internacional (ISO 2533) a $T_0 = 273.15\text{ K}$ ($0^\circ\text{C}$) y $P_0 = 101325\text{ Pa}$:
\begin{equation*}
c_s = 331.300000000000 \text{ m/s} \quad \text{(A } 0^\circ\text{C)}
\end{equation*}
A $20^\circ\text{C}$ ($293.15\text{ K}$):
\begin{equation*}
c_s = 343.210000000000 \text{ m/s}
\end{equation*}
\textbf{Incertidumbre / Margen de Error:} En aire real, la humedad relativa, la concentración de $\text{CO}_2$ y pequeños cambios térmicos introducen un margen de error experimental empírico de $\pm 0.1 \text{ m/s}$ ($\approx 0.03\%$).

\vspace{0.8cm}
\hrule
\vspace{0.8cm}

\section{Intensidad Sonora de Referencia Umbral ($I_0$)}

\subsection*{1. Significado, Símbolo y Explicación}
Representada como $I_0$, es el valor estándar internacional para el umbral mínimo de audición humana para una onda senoidal pura a la frecuencia de $1000 \text{ Hz}$. Es la constante de normalización utilizada para calcular el Nivel de Presión Sonora (SPL) en decibelios ($dB_{SPL}$):
\begin{equation*}
L_I = 10 \log_{10} \left( \frac{I}{I_0} \right)
\end{equation*}

\subsection*{2. Valores Típicos en Ingeniería y Ciencia}
\begin{equation*}
I_0 = 10^{-12} \text{ W/m}^2 \quad \text{o} \quad 1 \text{ pW/m}^2
\end{equation*}

\subsection*{3. Decimales Completos y Definición Exacta}
Fijada exactamente por norma internacional (ISO 80000-8):
\begin{equation*}
I_0 = 1.000000000000000 \times 10^{-12} \text{ W/m}^2
\end{equation*}
Presión acústica de referencia correspondiente en aire ($p_0$):
\begin{equation*}
p_0 = 2.000000000000000 \times 10^{-5} \text{ Pa} \quad (20 \text{ }\mu\text{Pa})
\end{equation*}
\textbf{Incertidumbre / Margen de Error:} Exacto por definición de referencia internacional de ingeniería ($u_r = 0$).

\end{document}
